\section{Related Work}
\label{sec:related}

\noindent\textbf{Event-based 3D reconstruction with NeRF.}
EventNeRF~\cite{eventnerf} introduced 3D-consistent photorealistic novel-view synthesis from a single colour event stream, training a NeRF in a self-supervised manner. Dynamic EventNeRF~\cite{dynamiceventnerf} extends this to general dynamic scenes from sparse multi-view event cameras plus sparse RGB frames, using cross-faded time-conditioned NeRFs. Deformable NeRF with RGB and event cameras~\cite{defnerfevent} models non-rigid deformation by fusing the asynchronous event stream with calibrated sparse RGB frames. These NeRF-based approaches are limited by slow training and rendering, and they reconstruct the \emph{full} scene from events, which bounds fidelity because events lack absolute intensity.

\noindent\textbf{Event-based 3D Gaussian Splatting.}
A recent and rapidly growing line replaces the implicit NeRF with explicit 3D Gaussians. Event-3DGS~\cite{event3dgs} is the first event-based 3DGS, jointly optimizing scene and sensor parameters with a high-pass-filter photovoltage module for noise robustness. Event3DGS~\cite{xiong2025event3dgs} targets high-speed robot egomotion, reconstructing high-fidelity structure purely from events with 95\% reduced compute. EventSplat~\cite{eventsplat} uses an event-to-video model for initialization and spline-interpolated poses to handle fast camera motion. EvGGS~\cite{evggs} is the first \emph{generalizable} event-based 3DGS, reconstructing scenes as Gaussians from only event input in a feed-forward manner. E2GS~\cite{e2gs} and EF-3DGS~\cite{ef3dgs} incorporate events into Gaussian Splatting to enhance rendering under challenging conditions. DiET-GS~\cite{dietgs} constrains 3DGS with an event double integral and a diffusion prior for motion deblurring. Event-Boosted Deformable 3DGS~\cite{eventboostdeform3dgs} uses a dynamic-static decomposition. EdMCGS~\cite{edmcgs} models scene motion as an event-driven Markov chain from extreme-low-frame-rate RGB and events, keeping the event transition active at inference. FastEventDGS~\cite{fasteventdgs} achieves 4D reconstruction from a single event camera with a continuous-trajectory parameterization and a local-patch event-motion loss. ERF-GS~\cite{erfgs} fuses events into the optimization and densification stages for fast-motion reconstruction from disjoint event-RGB viewpoints. \emph{All of these methods reconstruct the full scene from events and assume a moving camera to generate parallax; none exploit a known baseline scene or a static camera for dynamic-increment modeling, which is our setting.}

\noindent\textbf{Pose-free event-based reconstruction.}
IncEventGS~\cite{inceventgs} removes the known-pose assumption via a tracking-and-mapping SLAM pipeline. EvTrajGS~\cite{evtrajgs} enables joint pose-scene optimization from coarse pose priors. These works address the orthogonal problem of unknown camera pose; our method assumes a known static pose, which makes the pixel-to-Gaussian correspondence precomputable and is the key to high-temporal-resolution updates.

\noindent\textbf{Event camera calibration and pose estimation.}
Calibrating an event camera and estimating its pose relative to a 3D scene is a prerequisite for our known-pose assumption. Geometric frameworks recover absolute 6-DoF pose and velocity from 3D-line--event correspondences~\cite{liu2026eventpose}. Line-based methods track object pose by matching 2D event-lines to 3D models~\cite{liu2024lopet}. Active-marker systems use blinking LEDs with minimal 2-point solvers~\cite{yuan2026p2p}. These techniques provide the known camera pose our method relies on; combining them with EventGS is a promising direction for fully real-world deployment.

\noindent\textbf{Aperture problem and Contrast Maximimization.}
The aperture problem---that a local motion measurement only constrains the component perpendicular to the image edge---is classical in optical flow~\cite{barron1994aperture}. For event cameras, Gallego \etal~\cite{gallego2018contrast} introduced the Contrast Maximimization (CMax) framework, which estimates motion by maximizing the sharpness of the Image of Warped Events. Normal-flow methods~\cite{ren2024evlinear} formalize the partial-observability bias of event-based flow and provide linear solvers that account for it. We extend these ideas to the 3D-Gaussian setting: our per-Gaussian observability score is the structure-tensor quantity underlying Lucas-Kanade and Harris, and our CMax regularizer on the predicted 2D Gaussian flow is the 3D-Gaussian analogue of the classical event-warping objective.

\noindent\textbf{Event-to-video and event simulation.}
E2VID~\cite{rebecq2019e2vid} and follow-ups reconstruct intensity frames from events by inverting the event-generation model. A naive pipeline (events~$\to$~E2VID frames~$\to$~GS update) would inherit E2VID's frame-rate bottleneck and error accumulation, negating the microsecond resolution of events; our approach lets events act \emph{directly} on 3D Gaussian parameters. For experiments with real-world video, the v2e simulator~\cite{hu2021v2e} generates realistic DVS events from intensity frames; we use it as the event-generation backend for our real-data pipeline.

\noindent\textbf{3D Gaussian Splatting and dynamic extensions.}
3DGS~\cite{kerbl3Dgaussians} represents scenes as anisotropic Gaussians with differentiable rasterization, enabling real-time rendering and fast optimization. Dynamic extensions include Deformable 3D Gaussians~\cite{yang2023deformable3dgs} and 4D Gaussian Splatting~\cite{wu20244dgs}, which model time-varying content via deformation fields or 4D primitives. We build on the polynomial-motion formulation of StreamSplat~\cite{streamsplat2025}, which predicts per-Gaussian polynomial coefficients for online dynamic reconstruction from video streams; we re-purpose this decoder for event-conditioned increments.
